% Statement and cited-input packet. Target proof is separately embargoed.
% Frozen II source SHA256: 74497bfe327fae25c9cdcb17334b3bf42a9793a9df7c0a4404d97568f2f760f0
% Target owner: cdp:st:E14; tag: 01GL

\begin{proposition}[The Liu--Rollenske family and its transverse deformation]\hypertarget{tag-01GL}{}\label{cdp:st:E14}


Every member $Y(u)$ of the Liu--Rollenske family, with its scale
sufficiently small and nonzero as in its construction, satisfies
\[
 h^\bullet(Y(u),T_{Y(u)})=(1,2,2,0).
\]
Its Kuranishi germ is a smooth reduced surface. One parameter is
the given scale; a second parameter replaces its cuspidal fibre by
two rank-one nodal fibres. All sufficiently small deformations
retain a torus fibration and have algebraic dimension one.
\end{proposition}

% Cited input cdp:st:E1; statement or labelled formula only.
\begin{proposition}[Vertical divisors and the field of meromorphic functions]\hypertarget{tag-01GG}{}\label{cdp:st:E1}


Let $Y$ be a connected compact complex manifold of dimension $n\geq2$
with $H^2(Y,\mathbb R)=0$. Let $g:Y\to C$ be a surjective holomorphic
map to a connected compact Riemann surface, with connected fibres.
Suppose that its smooth fibres over a nonempty open subset of $C$ are
K\"ahler. Every effective divisor of $Y$ is supported on finitely many
fibres, and pullback induces an isomorphism of fields
\[
 g^*: \mathcal M(C)\xrightarrow{\ \sim\ }\mathcal M(Y).
\]
In particular $a(Y)=1$. Reducedness of the other fibres is unnecessary.
\end{proposition}

% Cited input cdp:st:E13; statement or labelled formula only.
\begin{proposition}[The Kuranishi germ of the two-parameter family]\hypertarget{tag-01GK}{}\label{cdp:st:E13}


For every member of the family of
\cite[Proposition labelled \texttt{family}]{Viaclovsky2026Two}, on the parameter
bidiscs with $\mu$ sufficiently close to zero and nonzero sufficiently
small linearization scale $\lambda$, one has
\[
 (h^0,h^1,h^2,h^3)(X_{\mu,\lambda},T_X)=(1,2,2,0).
\]
Its Kuranishi germ is a smooth reduced surface, and the displayed
$(\mu,\lambda)$-family is locally miniversal. In particular its
sufficiently small deformations all have algebraic dimension one.
\end{proposition}

% Cited input cdp:st:E6; statement or labelled formula only.
\begin{proposition}[Stability of a rank-two torus fibration]\hypertarget{tag-01GV}{}\label{cdp:st:E6}


Let $g:Y\to\mathbb P^1$ be as in \ref{cdp:st:E5} and suppose its
Kodaira--Spencer tensor has rank two at one regular value. For every
proper deformation $\mathcal Y\to(S,0)$ of $Y$ over a complex
analytic germ, including a nonreduced germ, there is, after shrinking
$S$, a holomorphic map
\[
 G:\mathcal Y\longrightarrow\mathbb P^1\times S
\]
over $S$ whose central restriction is $g$. Every $G_s$ is surjective
with connected fibres and smooth general fibre a complex two-torus.
With $b_2(Y)=0$, every sufficiently small deformation has algebraic
dimension one.
\end{proposition}

% Cited input cdp:st:E5; statement or labelled formula only.
\begin{proposition}[The infinitesimal map and the rank-two boundary]\hypertarget{tag-01GU}{}\label{cdp:st:E5}


Let $g:Y\to C$ be a proper surjective holomorphic map with connected fibres
from a smooth connected compact complex threefold to a nonsingular compact
connected complex curve. Its smooth general fibre is a two-dimensional complex
torus. There is a canonical exact sequence of vector spaces
\begin{equation}\hypertarget{tag-019Y}{}\label{cdp:eq:E-Leray-map}
0\to H^1(C,T_C)\xrightarrow{g^*}H^1(Y,g^*T_C)
 \to H^0(C,T_C\otimes R^1g_*\mathcal O_Y)\to0.
\end{equation}
On the smooth locus $U\subset C$, the relative tangent sequence
induces a sheaf morphism
\[
 \delta:T_C\otimes R^1g_*\mathcal O_Y|_U
                  \longrightarrow R^2g_*T_{Y/C}|_U.
\]
At a smooth torus fibre $F$, write its Kodaira--Spencer tensor as
$\kappa=u_1\otimes\eta_1+u_2\otimes\eta_2$, where
$\eta_1,\eta_2$ is a basis of $H^1(F,\mathcal O_F)$ and
$u_i\in\operatorname{Lie}(F)$. Then
\[
 \delta(a\eta_1+b\eta_2)
   =(bu_1-au_2)\otimes(\eta_1\wedge\eta_2).
\]
Its rank is the rank of $\kappa$. When this rank is two at one
regular value, the image of $H^1(Y,T_Y)\to H^1(Y,g^*T_C)$ lies
in $g^*H^1(C,T_C)$. For $C=\mathbb P^1$ this image is zero.
\end{proposition}

